\documentclass[10pt,a4paper]{amsart}
\usepackage{amsmath,amssymb,amsthm,graphicx}
% For commutative diagrams we use Xy-pic
\usepackage[all]{xy}

% We use 2cell for 2-commutative diagrams.
\xyoption{2cell}
% Theorem environments.
%
\theoremstyle{plain}
\newtheorem{theorem}[subsection]{Theorem}
\newtheorem{proposition}[subsection]{Proposition}
\newtheorem{lemma}[subsection]{Lemma}

\theoremstyle{definition}
\newtheorem{definition}[subsection]{Definition}
\newtheorem{example}[subsection]{Example}
\newtheorem{exercise}[subsection]{Exercise}
\newtheorem{situation}[subsection]{Situation}

\theoremstyle{remark}
\newtheorem{remark}[subsection]{Remark}
\newtheorem{remarks}[subsection]{Remarks}

\numberwithin{equation}{subsection}

% Macros
%
\def\lim{\mathop{\mathrm{lim}}\nolimits}
\def\colim{\mathop{\mathrm{colim}}\nolimits}
\def\Spec{\mathop{\mathrm{Spec}}}
\def\Hom{\mathop{\mathrm{Hom}}\nolimits}
\def\Ext{\mathop{\mathrm{Ext}}\nolimits}
\def\SheafHom{\mathop{\mathcal{H}\!\mathit{om}}\nolimits}
\def\SheafExt{\mathop{\mathcal{E}\!\mathit{xt}}\nolimits}
\def\Sch{\mathit{Sch}}
\def\Mor{\mathop{\mathrm{Mor}}\nolimits}
\def\Ob{\mathop{\mathrm{Ob}}\nolimits}
\def\Sh{\mathop{\mathit{Sh}}\nolimits}
\def\NL{\mathop{N\!L}\nolimits}
\def\CH{\mathop{\mathrm{CH}}\nolimits}
\def\proetale{{pro\text{-}\acute{e}tale}}
\def\etale{{\acute{e}tale}}
\def\QCoh{\mathit{QCoh}}
\def\Ker{\mathop{\mathrm{Ker}}}
\def\Im{\mathop{\mathrm{Im}}}
\def\Coker{\mathop{\mathrm{Coker}}}
\def\Coim{\mathop{\mathrm{Coim}}}

% Boxtimes
%
\DeclareMathSymbol{\boxtimes}{\mathbin}{AMSa}{"02}

%
% Macros for moduli stacks/spaces
%
\def\QCohstack{\mathcal{QC}\!\mathit{oh}}
\def\Cohstack{\mathcal{C}\!\mathit{oh}}
\def\Spacesstack{\mathcal{S}\!\mathit{paces}}
\def\Quotfunctor{\mathrm{Quot}}
\def\Hilbfunctor{\mathrm{Hilb}}
\def\Curvesstack{\mathcal{C}\!\mathit{urves}}
\def\Polarizedstack{\mathcal{P}\!\mathit{olarized}}
\def\Complexesstack{\mathcal{C}\!\mathit{omplexes}}
% \Pic is the operator that assigns to X its picard group, usage \Pic(X)
% \Picardstack_{X/B} denotes the Picard stack of X over B
% \Picardfunctor_{X/B} denotes the Picard functor of X over B
\def\Pic{\mathop{\mathrm{Pic}}\nolimits}
\def\Picardstack{\mathcal{P}\!\mathit{ic}}
\def\Picardfunctor{\mathrm{Pic}}
\def\Deformationcategory{\mathcal{D}\!\mathit{ef}}
\usepackage{cleveref}
\crefname{originallemma}{original lemma}{original lemmas}
\Crefname{originallemma}{Original lemma}{Original lemmas}
\makeatletter
\newlabel{lemma-nonvanishing-inherited}{{0DV7}{source}{}{sourcepointer-1}{}}
\newlabel{lemma-nonvanishing-inherited@cref}{{[originallemma][0][]0DV7}{[1][0][]source}{}{}{}}
\makeatother
\crefname{originallemma}{original lemma}{original lemmas}
\Crefname{originallemma}{Original lemma}{Original lemmas}
\makeatletter
\newlabel{lemma-equivalence-abelian-sheaves-point}{{04JQ}{source}{}{sourcepointer-2}{}}
\newlabel{lemma-equivalence-abelian-sheaves-point@cref}{{[originallemma][0][]04JQ}{[1][0][]source}{}{}{}}
\makeatother
\crefname{originaltheorem}{original theorem}{original theorems}
\Crefname{originaltheorem}{Original theorem}{Original theorems}
\makeatletter
\newlabel{theorem-colimit}{{09YQ}{source}{}{sourcepointer-3}{}}
\newlabel{theorem-colimit@cref}{{[originaltheorem][0][]09YQ}{[1][0][]source}{}{}{}}
\makeatother
\title{Reduction of the vanishing of l-power torsion cohomology to the maximal pro-l extension}
\author{The Stacks Project Authors}
\date{stacks-project:etale-cohomology.tex:local-snapshot-20261011; GNU FDL 1.2 or later}
\begin{document}
\maketitle
\noindent\textit{Editorial note.} All delivered mathematical text is the author's own, copied from the pinned Stacks Project chapter file etale-cohomology.tex; no statement or proof step was rewritten, reordered or removed.\par\medskip
\noindent\textit{Editorial note.} The proof is the author's complete argument for the equivalence, in both directions, and the two paragraphs after the set-up correspond exactly to those two directions.\par\medskip
\noindent\textit{Editorial note.} Every labelled reference occurring inside the counted spans is carried in assembly.external\_references as a non-load-bearing navigation pointer (same-chapter ones at their definition site, cross-chapter ones at the invocation inside the counted span), so no load-bearing mathematics is replaced by a stub. No bibliography citation occurs inside the counted spans.\par\medskip
\noindent\textit{Editorial note.} The checked chapter list for same-name collisions is the set of stacks chapters staged in this repository: chow, cohomology, divisors, etale-cohomology, flat, more-algebra, more-morphisms, morphisms, spaces-more-morphisms, varieties together with this card's chapter etale-cohomology.\par\medskip
\noindent\textit{Editorial note.} The preamble is the chapter's own native preamble excerpt already used for this source: the Xy-pic setup (native preamble lines 14-18) and the native theorem declarations with the mathematical macros (native preamble lines 154-217).\par\medskip
\noindent\textit{Allowed standard prerequisites (editorial).}\begin{itemize}
\item Maximal pro-l quotient of the absolute Galois group and the corresponding extension. Same chapter, the maximal pro-l subgroup H of G = Gal(K\textasciicircum{}sep/K) and the field L of the corresponding extension. Supplies the field L and the group whose cohomology is compared with that of G.
\item Comparison of Galois cohomology with etale cohomology of the spectrum of a field. Same chapter, the identification of H\textasciicircum{}n\_etale(Spec(K), F) with the Galois cohomology H\textasciicircum{}n(G, F\_sbar) for an l-power torsion sheaf. Supplies the translation used in both directions of the proof.
\item Spectral sequence or inflation-restriction argument for a closed subgroup of finite index. Same chapter, the behaviour of Galois cohomology under the passage to a pro-l subgroup and the vanishing of the cohomology of an l-group with l-power torsion coefficients. Supplies the argument that vanishing on the subgroup implies vanishing for the whole group in the coefficients considered.
\end{itemize}\medskip
\ifdefined\hypertarget\hypertarget{sourcepointer-1}{}\fi
\noindent\textit{Source reference.} Original lemma 0DV7: sources/stacks-etale-cohomology/pinned-original-source.tex, source lines 10543--10543. Navigation pointer invoked at proof line 10613; the counted span contains the pointer only and the cited result is not reproduced inside it, so the pointer is registered as a non-load-bearing navigation stub in the established form for Stacks material. No bibliography citation is involved.\par\medskip
\ifdefined\hypertarget\hypertarget{sourcepointer-2}{}\fi
\noindent\textit{Source reference.} Original lemma 04JQ: sources/stacks-etale-cohomology/pinned-original-source.tex, source lines 9078--9078. Navigation pointer invoked at proof line 10623; the counted span contains the pointer only and the cited result is not reproduced inside it, so the pointer is registered as a non-load-bearing navigation stub in the established form for Stacks material. No bibliography citation is involved.\par\medskip
\ifdefined\hypertarget\hypertarget{sourcepointer-3}{}\fi
\noindent\textit{Source reference.} Original theorem 09YQ: sources/stacks-etale-cohomology/pinned-original-source.tex, source lines 7400--7400. Navigation pointer invoked at proof line 10567; the counted span contains the pointer only and the cited result is not reproduced inside it, so the pointer is registered as a non-load-bearing navigation stub in the established form for Stacks material. No bibliography citation is involved.\par\medskip
\begin{lemma}
\label{lemma-reduce-to-l-group}
Let $\ell$ be a prime number and $n$ an integer $> 0$.
Let $K$ be a field with $G = Gal(K^{sep}/K)$ and let
$H \subset G$ be a maximal pro-$\ell$ subgroup with $L/K$
being the corresponding field extension. Then
$H^n_\etale(\Spec(K), \mathcal{F}) = 0$ for all
$\ell$-power torsion $\mathcal{F}$ if and only if
$H^n_\etale(\Spec(L), \underline{\mathbf{Z}/\ell\mathbf{Z}}) = 0$.
\end{lemma}

\begin{proof}
Write $L = \bigcup L_i$ as the union of its finite subextensions over $K$.
Our choice of $H$ implies that $[L_i : K]$ is prime to $\ell$.
Thus $\Spec(L) = \lim_{i \in I} \Spec(L_i)$ as in
Lemma \ref{lemma-nonvanishing-inherited}.
Thus we may replace $K$ by $L$ and assume that
the absolute Galois group $G$ of $K$ is a
profinite pro-$\ell$ group.

\medskip\noindent
Assume $H^n(\Spec(K), \underline{\mathbf{Z}/\ell\mathbf{Z}}) = 0$.
Let $\mathcal{F}$ be an $\ell$-power torsion sheaf on $\Spec(K)_\etale$.
We will show that $H^n_\etale(\Spec(K), \mathcal{F}) = 0$.
By the correspondence specified in
Lemma \ref{lemma-equivalence-abelian-sheaves-point}
our sheaf $\mathcal{F}$ corresponds to an $\ell$-power torsion
$G$-module $M$. Any finite set of elements $x_1, \ldots, x_m \in M$
must be fixed by an open subgroup $U$ by continuity.
Let $M'$ be the module spanned by the orbits of $x_1, \ldots, x_m$.
This is a finite abelian $\ell$-group
as each $x_i$ is killed by a power of $\ell$
and the orbits are finite. Since $M$ is the filtered colimit of
these submodules $M'$, we see that $\mathcal{F}$ is the filtered
colimit of the corresponding subsheaves $\mathcal{F}' \subset \mathcal{F}$.
Applying Theorem \ref{theorem-colimit} to this colimit, we reduce
to the case where $\mathcal{F}$ is a finite locally constant sheaf.

\medskip\noindent
Let $M$ be a finite abelian $\ell$-group with a continuous action
of the profinite pro-$\ell$ group $G$. Then there is a $G$-invariant
filtration
$$
0 = M_0 \subset M_1 \subset \ldots \subset M_r = M
$$
such that $M_{i + 1}/M_i \cong \mathbf{Z}/\ell \mathbf{Z}$ with
trivial $G$-action (this is a simple lemma on representation
theory of finite groups; insert future reference here).
Thus the corresponding sheaf $\mathcal{F}$ has a filtration
$$
0 = \mathcal{F}_0 \subset \mathcal{F}_1 \subset \ldots \subset
\mathcal{F}_r = \mathcal{F}
$$
with successive quotients isomorphic to
$\underline{\mathbf{Z}/\ell \mathbf{Z}}$.
Thus by induction and the long exact cohomology
sequence we conclude.
\end{proof}

\end{document}
